% SPDX-License-Identifier: Apache-2.0
% covers: Mesh
\datasheet{Mesh}{A lattice of network nodes as one unit, which lowers.}

\dsfacts{\code{lib/parts/src/bus/noc/mesh.rs}}{\code{txhdl\_parts::bus::noc::mesh::Mesh}}%
{\code{//docs:noc}}

\subsection*{Function}

\code{Mesh<W, H, N, ..>} is \code{N} nodes, \code{W} wide and \code{H}
high, in row-major order, held as one field, \code{nodes:
Units<Node, N>} (issue 635). Node \code{i} is at column \code{i \% W}
and row \code{i / W}, which the mesh gives it by tying its \code{col}
and \code{row} inputs to those constants.

The mesh's ports are the nodes' exits, four arrays of \code{N}:
\code{qx\_in} the requests into each node, \code{px\_in} the responses
into each, and \code{qx\_out} and \code{px\_out} what leaves by each.
Everything between the nodes is inside. The links are one array of
channels per direction of travel and per virtual channel. Node
\code{i}'s output to a direction is channel \code{i} of that
direction's array, and its neighbour reads that channel on the input
from the opposite side.

\subsection*{Parameters}

\begin{center}\footnotesize
\begin{tabular}{@{}lp{0.7\textwidth}@{}}
\toprule
Parameter & Meaning \\
\midrule
\code{W} & The width of the lattice, in nodes. \\
\code{H} & The height of the lattice, in nodes. \\
\code{N} & The number of nodes, which must be \code{W * H}. \\
\code{XB} & The width of a column coordinate, in bits: at least enough
for \code{W - 1}. \\
\code{YB} & The width of a row coordinate, in bits: at least enough
for \code{H - 1}. \\
\code{A} & The address width of the AXI link. \\
\code{D} & The data width of the AXI link. \\
\code{S} & The strobe width, which is \code{D / 8}. \\
\code{I} & The identifier width of the AXI link. \\
\bottomrule
\end{tabular}
\end{center}

The tables below are for \code{Mesh<3, 2, 6, 2, 1, 8, 8, 1, 2>}, the
mesh \code{ex\_mesh} makes: three nodes wide and two high, on a link
of eight-bit addresses and data.

\dstables{Mesh}

\subsection*{Behaviour}

\begin{itemize}
\item The mesh has no logic of its own and no register. What it holds
is its nodes' state, six bits each.
\item Node \code{i} reads from the north what the node above sent
down, from the south what the node below sent up, from the west what
the node to its west sent east, and from the east what the node to its
east sent west.
\item A link at an edge wraps to the other end of its row or column,
so that every channel has one driver and one reader. Routing is
dimension order, so no packet for a place on the lattice ever takes
one: nothing leaves the east edge eastward, since no column is east of
it.
\item \code{Default} refuses a mesh whose \code{N} is not
\code{W * H}, with a message naming both, since Rust cannot yet
compute \code{N} from \code{W} and \code{H}.
\end{itemize}

\subsection*{Verification}

\begin{itemize}
\item \code{ex\_mesh} runs \code{Mesh<3, 2, 6, ..>}. Every node sends a
request to each of the other five, and a response to the node at the
opposite corner. The run asserts that each packet leaves by the exit
of the node it was addressed to, and that all thirty requests and six
responses arrive.
\item The netlist, six node instances and the channels between them,
is checked against that run under nvc and Verilator, by
\code{//docs:mesh\_sim\_mesh\_tb\_test} and \code{//docs:mesh\_vsim\_test},
at every port of the mesh.
\end{itemize}

\subsection*{Used by}

\code{ex\_mesh}. The tests of the network and \code{//soc} wire their
lattices with \code{mesh::lattice}, which does not lower.

\subsection*{Limits}

Nothing refuses a packet addressed to a place off the lattice, a
column of \code{W} or more say: the wrapped links carry nothing only
while every destination is on it. The nodes' registers are checked
at the mesh's ports only, not inside it. Each node's switches move
one packet per cycle, not one per port.
